\documentclass[10pt]{article}

\usepackage[margin=1in]{geometry}
\usepackage{amsmath}
\usepackage{graphicx}
\usepackage{hyperref}
\usepackage{url}

\title{CURE-OR++: Object Recognition Robustness Under Native CURE-OR Challenges and Digital Transfer Pipelines}
\author{CURE-OR++ Project}
\date{\today}

\begin{document}

\maketitle

\begin{abstract}
Object-recognition systems are often evaluated on clean or synthetically corrupted images, but practical image use also includes acquisition artifacts, digital transfer, recompression, screenshots, and channel changes. CURE-OR++ builds a compact, reproducible benchmark layer around CURE-OR for measuring how vision models fail under object-centric native challenge conditions and planned real transfer pipelines. On the Full-CURE-OR v0.4 probe, eight usable baselines cover CLIP/OpenCLIP zero-shot models and frozen-feature prototype classifiers. The strongest current level-5 baseline, DINOv2 ViT-S/14 Prototype, reaches only 0.2766 mean level-5 accuracy despite 0.7520 clean accuracy. Across all eight usable baselines, grayscale salt-and-pepper noise, salt-and-pepper noise, and grayscale gaussian blur collapse to near-chance accuracy, with all models at the floor threshold. Confidence analysis shows that stronger zero-shot models can remain substantially overconfident under severe native challenges: OpenCLIP ViT-B/16 DataComp XL reaches 0.1451 level-5 accuracy while retaining 0.4997 mean confidence. A type-10 grayscale control confirms that grayscale conversion alone is damaging but does not explain the full native level-5 collapse. The remaining validation step is a 180-output real-transfer block covering messenger upload/download, phone screenshot/resave, and video-call frame capture pipelines.
\end{abstract}

\section{Introduction}

Robustness benchmarks often study isolated corruptions such as blur, noise, compression, or exposure shifts. Real image use is frequently more operational: images are captured, resized, recompressed, screenshotted, shared through applications, restored, or viewed through video-call pipelines. CURE-OR++ is designed around that narrower operational question: how do object-recognition baselines fail under reproducible object-centric challenge and transfer conditions?

This draft focuses on the Full-CURE-OR v0.4 controlled probe. The current repository also contains a smaller CURE-OR++ v0.1 Kaggle-ready artifact and a prepared real-transfer v0.2 protocol. The latter is the main remaining external-validity gap before a final public paper.

\paragraph{Current contributions.}
The current artifact provides:
\begin{enumerate}
  \item a manifest-driven Full-CURE-OR v0.4 probe over clean rows and native challenge rows;
  \item a multi-family baseline comparison covering CLIP/OpenCLIP zero-shot models, a stronger DataComp XL OpenCLIP variant, and four non-CLIP prototype classifiers;
  \item a level-5 consensus failure analysis across eight usable baselines;
  \item confidence and calibration analysis for the usable CLIP/OpenCLIP-family zero-shot rows;
  \item a type-10 grayscale control and paired color-versus-grayscale challenge-family analysis;
  \item a prepared real-transfer v0.2 protocol for messenger, screenshot, and video-call pipelines.
\end{enumerate}

\section{Benchmark Design}

CURE-OR++ uses manifest-driven evaluation. Each row links a source image, output image, object label, challenge family, recipe, severity, and metadata. This keeps aggregate results traceable to well-defined image groups and makes the benchmark easy to audit.

The Full-CURE-OR v0.4 track uses 500 clean probe rows and 38,999 native challenge probe rows across official challenge types 02-09 and 11-18. It also includes a type-10 grayscale no-challenge control. The real-transfer v0.2 track is prepared but not evaluated: it pins 30 clean mini-CURE-OR sources, three pipelines, two repeats per source and pipeline, and 180 expected transferred outputs.

\section{Models}

The current main table excludes the SigLIP Base P16 224 diagnostic row because the present zero-shot prompt protocol produced a clean-accuracy failure. The eight usable baseline rows are DINOv2 ViT-S/14 Prototype, ConvNeXt-Tiny Prototype, OpenCLIP ViT-B/16 DataComp XL, HGNetV2-B0 Prototype, CLIP ViT-B/16, MobileNetV3-Small Prototype, OpenCLIP ViT-B/32 LAION2B, and CLIP ViT-B/32.

\section{Results}

The paper-ready tables below are generated from tracked aggregate CSV files by
\path{scripts/build_paper_tables.py}. Native means are weighted by row count.

{\small
\begin{table}[t]
\centering
\caption{Full-CURE-OR v0.4 model leaderboard.}
\label{tab:full-cure-or-model-leaderboard}
\begin{tabular}{lllll}
\hline
Model & Clean & Native mean & Level 5 & Worst L5 \\
\hline
DINOv2 ViT-S/14 Prototype & 0.752 & 0.439 & 0.277 & Grayscale salt \& pepper noise (0.008) \\
ConvNeXt-Tiny Prototype & 0.616 & 0.344 & 0.224 & Grayscale salt \& pepper noise (0.008) \\
OpenCLIP ViT-B/16 DataComp XL & 0.546 & 0.256 & 0.145 & Salt \& pepper noise (0.008) \\
HGNetV2-B0 Prototype & 0.624 & 0.255 & 0.122 & Gaussian blur (0.010) \\
CLIP ViT-B/16 & 0.444 & 0.193 & 0.099 & Salt \& pepper noise (0.010) \\
MobileNetV3-Small Prototype & 0.556 & 0.194 & 0.096 & Gaussian blur (0.010) \\
OpenCLIP ViT-B/32 LAION2B & 0.412 & 0.170 & 0.089 & Grayscale salt \& pepper noise (0.008) \\
CLIP ViT-B/32 & 0.350 & 0.153 & 0.074 & Salt \& pepper noise (0.006) \\
\hline
\end{tabular}
\end{table}

\begin{table}[t]
\centering
\caption{Consensus hardest level-5 native CURE-OR challenges.}
\label{tab:full-cure-or-consensus-failures}
\begin{tabular}{lllll}
\hline
Rank & Type & Challenge & Mean acc. & Floor \\
\hline
1 & 18 & Grayscale salt \& pepper noise & 0.009 & 8/8 \\
2 & 09 & Salt \& pepper noise & 0.010 & 8/8 \\
3 & 14 & Grayscale gaussian blur & 0.011 & 8/8 \\
4 & 05 & Gaussian blur & 0.015 & 7/8 \\
5 & 16 & Grayscale dirty lens 1 & 0.026 & 5/8 \\
6 & 07 & Dirty lens 1 & 0.089 & 0/8 \\
\hline
\end{tabular}
\end{table}

\begin{table}[t]
\centering
\caption{Type-10 grayscale control versus native level-5 accuracy.}
\label{tab:full-cure-or-grayscale-control}
\begin{tabular}{llll}
\hline
Model & Gray control & Native L5 & Control - L5 \\
\hline
DINOv2 ViT-S/14 Prototype & 0.741 & 0.277 & 0.465 \\
ConvNeXt-Tiny Prototype & 0.491 & 0.224 & 0.267 \\
OpenCLIP ViT-B/16 DataComp XL & 0.391 & 0.145 & 0.246 \\
HGNetV2-B0 Prototype & 0.443 & 0.122 & 0.321 \\
CLIP ViT-B/16 & 0.317 & 0.099 & 0.217 \\
MobileNetV3-Small Prototype & 0.285 & 0.096 & 0.189 \\
OpenCLIP ViT-B/32 LAION2B & 0.246 & 0.089 & 0.157 \\
CLIP ViT-B/32 & 0.224 & 0.074 & 0.150 \\
\hline
\end{tabular}
\end{table}

}

\subsection{Level-5 Robustness}

The strongest current level-5 row is DINOv2 ViT-S/14 Prototype: 0.7520 clean accuracy, 0.4393 native mean accuracy, and 0.2766 native level-5 mean accuracy. The strongest zero-shot contrastive row is OpenCLIP ViT-B/16 DataComp XL: 0.5460 clean accuracy, 0.2561 native mean accuracy, and 0.1451 native level-5 mean accuracy. Stronger features and pretraining improve the leaderboard, but they do not remove the severe native challenge collapse.

\subsection{Consensus Failure Ordering}

The top three consensus level-5 failures are grayscale salt-and-pepper noise, salt-and-pepper noise, and grayscale gaussian blur. All eight usable baselines are at the floor threshold on these three challenges. Pairwise level-5 rank correlations range from 0.892 to 0.988, suggesting a stable severe-challenge hardness core rather than arbitrary per-model failure lists.

\subsection{Confidence and Calibration}

The confidence pass currently covers the usable CLIP/OpenCLIP-family zero-shot rows. OpenCLIP ViT-B/16 DataComp XL reaches 0.1451 level-5 accuracy while retaining 0.4997 mean confidence. Its worst level-5 overconfidence case is grayscale salt-and-pepper noise, with 0.0100 accuracy, 0.7393 mean confidence, and a 0.9760 high-confidence wrong rate. Severe native challenge failure can therefore preserve high confidence.

\subsection{Grayscale Control}

The type-10 grayscale no-challenge control is damaging but does not explain the full level-5 collapse. For OpenCLIP ViT-B/16 DataComp XL, grayscale control accuracy is 0.3908 while native level-5 mean accuracy is 0.1451. Paired channel-effect analysis also shows that every usable model is worse on grayscale level-5 native challenges than on paired color level-5 native challenges.

\section{Related Work}

CURE-OR++ should be positioned as a compact, object-centric robustness artifact rather than a replacement for broad VLM robustness suites. Nearby benchmarks include R-Bench~\cite{rbench}, MLLM-IC~\cite{mllmic}, VLM-RobustBench~\cite{vlmrobustbench}, and MMD-Bench/CLEAR~\cite{mmdbenchclear}. The intended differentiator is not scale, but transparent generation, paired metadata, small object-recognition scope, and failure tables that are easy to inspect by object class, challenge recipe, severity, and model family.

\section{Limitations}

Real-transfer v0.2 is prepared but not yet evaluated. Full-CURE-OR v0.4 is a controlled probe rather than an exhaustive evaluation of every source image. Prototype classifiers and zero-shot contrastive models are useful robustness contrasts, but they are not identical evaluation protocols. Confidence and calibration analysis currently covers the usable CLIP/OpenCLIP-family zero-shot rows rather than every prototype row.

\section{Real-Transfer Placeholder}

This section should be replaced after the 180-output real-transfer v0.2 block is collected, activated, and evaluated. The prepared command is:

\begin{verbatim}
.venv/bin/python scripts/activate_real_transfer_protocol.py --require-ready
\end{verbatim}

The planned real-transfer rows evaluate:

\begin{itemize}
  \item CLIP ViT-B/16;
  \item CLIP ViT-B/32;
  \item OpenCLIP ViT-B/32 LAION2B;
  \item OpenCLIP ViT-B/16 DataComp XL.
\end{itemize}

\section{Conclusion}

CURE-OR++ currently provides a strong pre-paper benchmark artifact. The Full-CURE-OR v0.4 probe shows severe native challenge failures across model families, consensus floor failures under salt-and-pepper and blur variants, and confidence-preserving errors in stronger zero-shot models. The critical next step is real-transfer v0.2, which will test whether actual app and device transfer pipelines reproduce related failure behavior.

\bibliographystyle{plain}
\bibliography{references}

\end{document}
